% Personal letter template using ashletter.cls.
% This is written for my own reuse, so the comments below are practical reminders,
% not package documentation.

\documentclass[11pt]{ashletter}
\usepackage[
  letterpaper,
  left=0.7in,
  right=0.7in,
  top=1.0in,
  bottom=1.0in
]{geometry}
\usepackage{soul}
\usepackage{xcolor}

% Set the date explicitly.
% Common options:
% \date{\today}
% \date{April 15, 2026}
\date{\today}

% Subject line shown above the salutation.
\subject{REPLACE WITH SUBJECT LINE}

% Examples of header customization.
% Only change these if I actually want to override the default SBU/Kumar-lab header.
%
% \setheadertitle{College of Arts and Sciences}
% \setheaderdepartment{Department of Physics and Astronomy}
% \setheaderinstitution{Stony Brook University}
% \setheaderaddress{Stony Brook, NY 11794}
% \setheadername{Aishwarya Kumar}
% \setheaderrole{Assistant Professor of Physics \& Astronomy}
% \setheaderphone{8144418969}
% \setheaderemail{aishwarya.kumar@stonybrook.edu}
% \setheaderwebsite{https://sites.google.com/stonybrook.edu/kumarlab}
% \setheaderlogofile{SBU stack_2clr_rgb_300ppi.png}
%
% Useful spacing/layout knobs if the template needs small visual adjustments.
% Margins are set here in the document, not hidden in the class.
% The header follows the document margins automatically.
% By default, the divider position is derived from the rendered logo width plus
% a small right-side pad, so I usually should not need to tune it.
%
% Header top offset usually needs to match the negative of the top margin if I want
% the header flush to the top of the page.
% \setheadertopoffset{-1.0in}
% \setheaderleftwidth{3.9in}          % manual override if I explicitly want one
% \setheaderlogorightpad{0.18in}      % tweak divider distance from the logo
\setheaderrightxshift{0.2in}      % move only the right header block left/right
% \setheaderrightyshift{0.2in}       % move only the right header block down/up
% \setheadertextvshift{-0.48in}       % base internal vertical alignment of right header text
% \setheaderbottomgap{0.8\baselineskip}
% \setrecipientgap{0.8\baselineskip}
% \setsubjectgap{0.8\baselineskip}
% \setclosinggap{1.2\baselineskip}
% \setsignaturevspace{0.1in}
% \setheadertopoffset{-.8in}

\begin{document}

% Recipient block.
% Use \\ between lines. Keep this block compact.
\begin{letter}{%
RECIPIENT NAME\\
RECIPIENT TITLE / COMMITTEE / ORGANIZATION\\
RECIPIENT INSTITUTION OR ADDRESS
}

% Opening should just be the salutation.
\opening{Dear RECIPIENT,}

% Body notes to self:
% - Keep paragraphs short and readable.
% - If this is a support/recommendation letter, open with the recommendation immediately.
% - Avoid manual \vspace or \newpage unless absolutely necessary.
% - Let the class handle layout first; only tweak class spacing if the template genuinely needs it.

FIRST PARAGRAPH GOES HERE. Explain the purpose of the letter immediately.

SECOND PARAGRAPH GOES HERE. Add the most important supporting details.

THIRD PARAGRAPH GOES HERE. Keep going with concrete examples, evidence, or context.

FINAL PARAGRAPH GOES HERE. End with a clear recommendation / request / conclusion.



Sincerely,\\ \\

% Signature block.
% This is the stable manual API. Use SignatureRow + Signer.
%
% One signer:
% \SignatureRow{l}
% \Signer{AK_signature.png}{Aishwarya Kumar}
% \EndSignatureRow
%
% Two signers:
% \SignatureRow{ll}
% \Signer{sig1.png}{Name One} &
% \Signer{sig2.png}{Name Two}
% \EndSignatureRow
%
% Three signers:
% \SignatureRow{lll}
% \Signer{sig1.png}{Name One} &
% \Signer{sig2.png}{Name Two} &
% \Signer{sig3.png}{Name Three}
% \EndSignatureRow
%
% If I need more than three signers, just stack multiple SignatureRow blocks.
\SignatureRow{l}
\Signer{AK_signature.png}{Aishwarya Kumar}
\EndSignatureRow

\end{letter}
\end{document}
